% Generated draft: requires mathematical and source review.
\documentclass[11pt]{article}
\usepackage[margin=25mm]{geometry}
\usepackage{fontspec}
\setmainfont{Latin Modern Roman}
\usepackage{amsmath,amssymb,mathtools,unicode-math}
\setmathfont{Latin Modern Math}
\usepackage{enumitem,needspace,xurl,hyperref}
\setlength{\parindent}{0pt}
\setlength{\parskip}{5pt}
\setlength{\emergencystretch}{3em}
\newcommand{\N}{\mathbb{N}}
\newcommand{\Z}{\mathbb{Z}}
\newcommand{\Q}{\mathbb{Q}}
\newcommand{\R}{\mathbb{R}}
\newcommand{\C}{\mathbb{C}}
\newcommand{\F}{\mathbb{F}}
\newcommand{\A}{\mathbb{A}}
\newcommand{\PP}{\mathbb{P}}
\newcommand{\E}{\mathbb{E}}
\newcommand{\eps}{\varepsilon}
\newcommand{\dd}{\,\mathrm{d}}
\begin{document}
\section*{UnsolvedMath ID: 2100408}
\textbf{Source problem code:} catalogue entry
\subsection{Definitions and mathematical statement}
Let \(K_1\subset\operatorname{int}K_2\subset\R^d\), \(d\ge2\), be convex bodies with boundaries \(\Sigma_1,\Sigma_2\). To give the usual inner-billiard meaning, suppose these boundaries have enough smoothness and strict convexity to define the reflection maps away from tangencies. An oriented line has a unit direction \(v\). At its forward intersection \(x\) with \(\Sigma_i\), specular reflection replaces \(v\) by
\[
v'=v-2\langle v,\nu_i(x)\rangle\nu_i(x),
\]
where \(\nu_i(x)\) is the unit normal, producing the outgoing oriented line \(B_i(\ell)\). Let \(\mathcal L\) denote the intended set of oriented lines meeting both hypersurfaces, with tangencies treated according to a specified convention.

The supplied commuting-inner-billiards problem asks whether
\[
B_1B_2=B_2B_1\text{ on the intended common domain}
\quad\Longrightarrow\quad
\Sigma_1,\Sigma_2\text{ are confocal ellipsoids}.
\]
The conclusion means that, after one Euclidean change of coordinates, there are real parameters \(a_1,\ldots,a_d,\lambda_1,\lambda_2\), with \(a_j-\lambda_i>0\), such that
\[
\Sigma_i=\left\{x:\sum_{j=1}^d\frac{x_j^2}{a_j-\lambda_i}=1\right\},\qquad i=1,2.
\]
The premise must specify whether \(B_i\) preserve \(\mathcal L\) and commute as self-maps there, or whether equality is required only where both compositions are defined. The excerpt does not state this domain convention, differentiability assumptions, strict convexity, or treatment of tangencies. These are necessary qualifications of the reflection formulation, not silently added source hypotheses.

The supplied research summary reports the inner problem solved in the plane and higher dimensions with confocal-ellipsoid conclusion. It separately reports a planar \emph{outer}-billiard conclusion of concentric homothetic ellipses and says the higher-dimensional outer/dual version appears open. Those outer statements are contextual claims, not replacements for the source's inner-billiard target. No supporting result citations are supplied.
\subsection{Short English statement}
Must two nested convex hypersurfaces with commuting inner-billiard reflections be confocal ellipsoids? The supplied summary reports yes, but the map domains, boundary regularity, and exact supporting theorem need review.
\subsection{Sources}
\begin{itemize}
\item Local UnsolvedMath registry, permanent ID \(2100408\), exact source code \texttt{catalogue entry}, ``Open Problems in Integrable Systems --- Geometry of caustics, invariant surfaces, and commuting billiard maps.''
\item \emph{Open Problems, Questions, and Challenges in Finite-Dimensional Integrable Systems} (2018), Problem 4.8, attributed to A. Glutsyuk and S. Tabachnikov; \url{https://arxiv.org/abs/1804.03737}. Source-TeX extraction; recorded access 2026-07-29.
\item Supplied research summary reports resolutions for inner billiards and planar outer billiards, but provides no authoritative paper, theorem number, or bibliographic citation for those resolutions.
\end{itemize}
\end{document}
